\documentclass[11pt]{article}
\usepackage[margin=1in]{geometry}
\usepackage{amsmath,amssymb,amsthm}
\usepackage{hyperref}

\newtheorem{theorem}{Theorem}
\newtheorem{lemma}[theorem]{Lemma}
\newtheorem{proposition}[theorem]{Proposition}
\newtheorem{corollary}[theorem]{Corollary}
\newtheorem{definition}[theorem]{Definition}
\newtheorem{remark}[theorem]{Remark}

\title{The EML tree algebra: a Lean~4 formalization}
\author{}
\date{}

\begin{document}
\maketitle

\begin{abstract}
Odrzywołek~\cite{odrzywolek2026} showed that a single binary operator
$\operatorname{eml}(x,y) = e^x - \ln y$, together with the constant~$1$,
suffices to express every elementary function.
We formalize the resulting tree algebra in Lean~4 and develop its
equational theory.
Differentiation becomes a total, computable function on trees, defined
by a single recursive clause.
We prove, as machine-checked rewrite chains with no axioms and no
dependence on Mathlib, that:
(i)~the derivative of every variable-free tree rewrites to zero;
(ii)~$D(e^f) \to e^f \cdot D(f)$ and $D(\ln f) \to D(f)/f$ for
arbitrary subtrees~$f$;
(iii)~differentiation is linear (sum, difference, negation, and
constant-multiple rules).
We further observe that every domain restriction in elementary
mathematics---division by zero, $\ln 0$, indeterminate
powers---traces to a single point of partiality in the EML operator,
and prove that this partiality is the unique source of non-confluence
in the rewrite system: the expression $0 \cdot (1/0)$ rewrites to
both~$0$ and~$1$.
\end{abstract}

%% -----------------------------------------------------------------
\section{Background}
%% -----------------------------------------------------------------

We briefly recall the EML operator introduced in~\cite{odrzywolek2026}.

\begin{definition}[EML operator]
The \emph{Exp-Minus-Log} operator is
\begin{equation}\label{eq:eml}
  \operatorname{eml}(x,y) \;=\; e^{x} - \ln y.
\end{equation}
An \emph{EML tree} is a term generated by the grammar
$S \to 1 \mid \operatorname{eml}(S,S)$,
optionally extended with input variables $x, y, \ldots$
\end{definition}

As shown in~\cite{odrzywolek2026}, every standard elementary function
can be expressed as an EML tree.  The key derived operations are:
\begin{alignat}{2}
  e^z        &= \operatorname{eml}(z, 1), &\qquad
  \ln z      &= \operatorname{eml}(1, \operatorname{eml}(\operatorname{eml}(1,z),1)),
  \label{eq:exp-ln} \\
  x - y      &= \operatorname{eml}(\ln x,\; e^y), &\qquad
  {-z}       &= 0 - z, \quad\text{where } 0 = \ln 1,
  \label{eq:sub-neg} \\
  x + y      &= x - (-y), &\qquad
  1/z        &= e^{-\ln z},
  \label{eq:add-inv} \\
  x \times y &= e^{\ln x + \ln y}, &\qquad
  x / y      &= x \times (1/y).
  \label{eq:mul-div}
\end{alignat}
Constants such as $e = \operatorname{eml}(1,1)$, $\;i = e^{\ln(-1)/2}$,
and $\pi = -i\ln(-1)$ are likewise finite trees.
Trigonometric and hyperbolic functions follow from Euler's formula.

We take these constructions as given and refer the reader
to~\cite{odrzywolek2026} for the exhaustive search methodology and
completeness argument.  The present paper is concerned exclusively with
what can be \emph{proved} about the resulting tree algebra, using
machine-checked reasoning in Lean~4.

%% -----------------------------------------------------------------
\section{The rewrite system}
\label{sec:rewrite}
%% -----------------------------------------------------------------

\subsection{Rewrite rules}

We equip EML trees with a directed rewrite relation $\to$, generated by
10~equational axioms and 2~congruence rules.

\medskip\noindent
\textbf{Cancellation.}
\begin{align}
  e^{\ln z} &\to z, &
  \ln(e^z)  &\to z.
  \label{eq:cancel}
\end{align}

\noindent
\textbf{Field identities.}
\begin{align}
  a - a &\to 0, &
  a \cdot 1 &\to a, &
  a \cdot 0 &\to 0, &
  -(-a) &\to a.
  \label{eq:field}
\end{align}

\noindent
\textbf{Structure.}
\begin{align}
  e^{a+b}    &\to e^a \cdot e^b, &
  \ln(a \cdot b) &\to \ln a + \ln b, \label{eq:hom}\\
  (a+b)+c    &\to a+(b+c), &
  a + b      &\to b + a. \label{eq:add-struct}
\end{align}

\noindent
\textbf{Congruence.} If $a \to a'$ then
$\operatorname{eml}(a,b) \to \operatorname{eml}(a',b)$ and
$\operatorname{eml}(a,b) \to \operatorname{eml}(a,b')$
(rewriting inside either child).

\medskip
We write $a \to^{*} b$ for the reflexive-transitive closure of $\to$.
Several convenient derived rules follow from the axioms above:
\begin{itemize}
\item $e^0 \to 1$ is the instance $z = 1$ of $e^{\ln z} \to z$,
  since $0 = \ln 1$ by definition.
\item $a - 0 \to^{*} a$ follows from $e^0 \to 1$ and
  $e^{\ln a} \to a$ (two applications of cancellation).
\item $0 + a \to^{*} a$ and $a + 0 \to^{*} a$ follow from
  $a - a \to 0$, $a - 0 \to^{*} a$, and commutativity.
\item $1 \cdot a \to^{*} a$ and $0 \cdot a \to^{*} 0$ follow from
  the right-hand versions via commutativity of multiplication, which
  is itself derived from commutativity of addition
  applied inside $\exp$.
\item $(1/(1/a)) \to^{*} a$ follows from $\ln(e^z) \to z$,
  $-(-a) \to a$, and $e^{\ln z} \to z$.
\end{itemize}
The Lean formalization includes all of these as named lemmas for
proof engineering convenience (20 rules total), but only the
10~axioms above are independent.

\subsection{Congruence lifting}

A rewrite inside a subexpression propagates to the surrounding context.
Because EML's derived operations (addition, multiplication, etc.)
are defined as specific tree shapes, rewriting inside them requires
explicit \emph{congruence lemmas}: if $a \to^{*} a'$ then
$a + b \to^{*} a' + b$, etc.  We prove these for all derived operations
(15~lemmas), building on the two primitive congruence rules.

%% -----------------------------------------------------------------
\section{Differentiation}
\label{sec:diff}
%% -----------------------------------------------------------------

\subsection{Definition}

Differentiation is a total function $D_x : \texttt{Eml} \to \texttt{Eml}$
defined by structural recursion:
\begin{align}
  D_x(1) &= 0, \notag\\
  D_x(v) &= \begin{cases} 1 & \text{if } v = x, \\ 0 & \text{otherwise,} \end{cases} \notag\\
  D_x(\operatorname{eml}(A,B)) &=
    e^A \cdot D_x(A) \;-\; D_x(B) \cdot (1/B).
  \label{eq:diff}
\end{align}
This is the chain rule applied to~\eqref{eq:eml}, but it requires no
semantic interpretation---it is pure tree surgery.  The output is always
another EML tree; differentiation never fails or diverges.

\subsection{Constants vanish}

\begin{theorem}\label{thm:const-zero}
If $t$ is a variable-free EML tree, then $D_x(t) \to^{*} 0$.
\end{theorem}

\begin{proof}
By structural induction on~$t$.  The base case $D_x(1) = 0$ is
immediate.  For $t = \operatorname{eml}(A,B)$ with $A,B$ variable-free,
the induction hypothesis gives $D_x(A) \to^{*} 0$ and
$D_x(B) \to^{*} 0$.  Then
\[
  e^A \cdot D_x(A) \;\to^{*}\; e^A \cdot 0 \;\to^{*}\; 0,
  \qquad
  D_x(B) \cdot (1/B) \;\to^{*}\; 0 \cdot (1/B) \;\to^{*}\; 0,
\]
and $0 - 0 \to^{*} 0$.  The only rules used are zero absorption
($a \cdot 0 \to 0$, $0 \cdot a \to 0$) and self-cancellation
($a - a \to 0$).  No exp, ln, or numerical evaluation appears.
\end{proof}

This single theorem immediately yields $D(e) = 0$, $D(\pi) = 0$,
$D(i) = 0$, etc., for every EML-encoded constant.

\subsection{Chain rules}

\begin{theorem}\label{thm:chain-exp}
For any subtree $f$, \;$D_x(e^f) \to^{*} e^f \cdot D_x(f)$.
\end{theorem}

\begin{proof}
Since $e^f = \operatorname{eml}(f, 1)$, we have
$D_x(e^f) = e^f \cdot D_x(f) - D_x(1) \cdot (1/1)$.
The second term rewrites: $D_x(1) = 0$, so $0 \cdot (1/1) \to 0$,
and $(\cdots) - 0 \to e^f \cdot D_x(f)$.
\end{proof}

\begin{theorem}\label{thm:chain-ln}
For any subtree $f$, \;$D_x(\ln f) \to^{*} D_x(f) \cdot (1/f)$.
\end{theorem}

\begin{proof}
The proof chains through five intermediate forms.
Since $\ln f = \operatorname{eml}(1, e^{\operatorname{eml}(1,f)})$,
differentiating and simplifying:
\begin{enumerate}
\item $D(1) = 0$ kills the left term, leaving
  $-\bigl(D(e^{\operatorname{eml}(1,f)}) \cdot (1/e^{\operatorname{eml}(1,f)})\bigr)$.
\item Apply the exponential chain rule (Theorem~\ref{thm:chain-exp})
  to expand the derivative of $e^{\operatorname{eml}(1,f)}$.
\item Apply the node rule to expand
  $D(\operatorname{eml}(1,f)) \to^{*} -(D(f) \cdot (1/f))$.
\item Cancel $e^{\operatorname{eml}(1,f)}$ against its inverse using
  the algebraic lemma $(ab) \cdot (1/a) \to^{*} b$
  (Lemma~\ref{lem:cancel}).
\item Double negation: $-(-(D(f) \cdot (1/f))) \to^{*} D(f) \cdot (1/f)$.
\end{enumerate}
\end{proof}

\begin{lemma}[Left cancellation]\label{lem:cancel}
$(a \cdot b) \cdot (1/a) \;\to^{*}\; b$.
\end{lemma}

\begin{proof}
An eight-step rewrite chain working entirely inside the outermost
$\exp$, using $\ln(a \cdot b) \to \ln a + \ln b$, the cancellation
$\ln(e^{-\ln a}) \to -\ln a$, associativity and commutativity of
addition, $\ln a + (-\ln a) \to 0$, $0 + \ln b \to \ln b$, and
finally $e^{\ln b} \to b$.
\end{proof}

\subsection{Linearity}

Using both chain rules and the cancellation lemma, we derive:

\begin{theorem}[Linearity of differentiation]\label{thm:linear}
For any subtrees $a, b, f$ and any variable-free $c$:
\begin{enumerate}
\item $D(a - b) \to^{*} D(a) - D(b)$,
\item $D(-f) \to^{*} -(D(f))$,
\item $D(a + b) \to^{*} D(a) + D(b)$,
\item $D(c \cdot f) \to^{*} c \cdot D(f)$.
\end{enumerate}
\end{theorem}

\begin{proof}
(1) The subtraction rule follows from expanding $a - b =
\operatorname{eml}(\ln a, e^b)$ and simplifying each half: the left
via $e^{\ln a} \to a$ and the logarithmic chain rule, the right via
the exponential chain rule, each followed by cancellation.
(2) follows from (1) with $a = 0$ and Theorem~\ref{thm:const-zero}.
(3) follows from (1) and (2) since $a + b = a - (-b)$.
(4) requires expanding $c \cdot f = e^{\ln c + \ln f}$, applying (3)
to the sum of logarithms, using Theorem~\ref{thm:const-zero} to kill
$D(\ln c)$ (since $c$ is ground, $\ln c$ is ground), applying the
logarithmic chain rule, and cancelling.
\end{proof}

\subsection{Concrete results}

As immediate corollaries:
\begin{align}
  D_x(e^x) &\to^{*} e^x, \label{eq:dexp}\\
  D_x(\ln x) &\to^{*} 1/x. \label{eq:dln}
\end{align}

%% -----------------------------------------------------------------
\section{Algebraic identities}
\label{sec:algebra}
%% -----------------------------------------------------------------

The following identities are proved as rewrite chains:
\begin{align}
  a + (-a) &\to^{*} 0, &
  a \cdot (1/a) &\to^{*} 1, \label{eq:inv} \\
  (a \cdot b) \cdot (1/a) &\to^{*} b, &
  a \cdot ((1/a) \cdot b) &\to^{*} b, \label{eq:cancel2}
\end{align}
as well as commutativity and associativity of addition (as single
steps) and commutativity of multiplication (derived from commutativity
of addition on the logarithms).

Associativity of multiplication is proved by \emph{joinability}:
both $a \cdot (b \cdot c)$ and $(a \cdot b) \cdot c$ reduce to
$e^{\ln a + (\ln b + \ln c)}$ via $\ln(e^z) \to z$ and
$\ln(a \cdot b) \to \ln a + \ln b$, meeting at a common reduct
via the associativity rule for addition.

\begin{remark}
These identities are notable for what they do \emph{not} use.
The proofs manipulate tree structure using the 20 rewrite rules; no
evaluation function, no real numbers, no field axioms are invoked.
The ``field-like'' behavior emerges from the tree encodings
of $+, \times, -, 1/x$ and the interaction of $\exp/\ln$
cancellation with associativity and commutativity of addition.
\end{remark}

%% -----------------------------------------------------------------
\section{Partiality and non-confluence}
\label{sec:partiality}
%% -----------------------------------------------------------------

\subsection{The single source of partiality}

Every domain restriction in elementary mathematics traces to $\ln$ at
zero.  The key observation is structural:
\begin{alignat}{2}
  1/z &= e^{-\ln z} &\qquad&\Longrightarrow\quad
    \text{$1/z$ undefined iff $\ln z$ undefined},
  \label{eq:inv-ln} \\
  z^w &= e^{w \ln z} &\qquad&\Longrightarrow\quad
    \text{$z^w$ undefined (for general $w$) iff $\ln z$ undefined}.
  \label{eq:pow-ln}
\end{alignat}
Division by zero is not an independent restriction---it is a
\emph{consequence} of $\ln 0$ being undefined:
\[
  1/0 = e^{-\ln 0}.
\]
Thus, all genuine partiality in elementary mathematics reduces to a
single predicate on a single point: $\ln(z)$ requires $z \neq 0$.

\begin{remark}
Some EML tree encodings pass through $\ln$ at points where the
encoded operation is actually total.  For example,
$a - b = \operatorname{eml}(\ln a, e^b)$ routes through $\ln a$,
even though $0 - b = -b$ is perfectly well-defined.  Similarly,
$a \times b = e^{\ln a + \ln b}$ passes through $\ln a$ and
$\ln b$, even though $0 \times b = 0$ is fine.
These are \emph{removable} singularities of the encoding: the rewrite
rules $0 \cdot a \to 0$ and $a - a \to 0$ give correct results
without ever evaluating $\ln 0$.  The essential, non-removable
singularities are exactly those of $1/z$ and $z^w$---both inherited
from $\ln$.
\end{remark}

\subsection{Non-confluence at division by zero}

\begin{theorem}\label{thm:nonconf}
$0 \cdot (1/0) \to^{*} 0$ \;and\; $0 \cdot (1/0) \to^{*} 1$.
\end{theorem}

\begin{proof}
For the first reduction: $0 \cdot a \to 0$ (zero absorption).

For the second: $0 \cdot (1/0) \to^{*} (1/0) \cdot 0$ (commutativity
of multiplication) $\to^{*} 1$ (multiplicative inverse,
$(1/a) \cdot a \to^{*} 1$).
\end{proof}

\begin{remark}
The two rewrite paths are each individually sound---zero absorbs under
multiplication, and inverses cancel.  They give contradictory results
only at $a = 0$, where $1/a$ is undefined.  The rewrite rules
$0 \cdot a \to 0$ and $a \cdot (1/a) \to 1$ are both necessary for
the algebraic theory; they are ``regularizations'' that give correct
answers at removable singularities.  At the essential singularity
$1/0$, the regularizations conflict.

If the rewrite system is confluent away from $1/0$---an open
conjecture---then this theorem identifies division by zero as the
\emph{unique} source of non-confluence.
\end{remark}

%% -----------------------------------------------------------------
\section{Formalization}
\label{sec:formalization}
%% -----------------------------------------------------------------

The Lean~4 formalization comprises approximately 1400 lines across
seven files.  All theorems described in this paper are fully
machine-checked with zero uses of \texttt{sorry} (Lean's axiom-deferral
mechanism) and no dependence on Mathlib or any external library.

The core files and their contents:

\medskip
\begin{tabular}{lrl}
\textbf{File} & \textbf{Lines} & \textbf{Contents} \\
\hline
\texttt{Basic.lean}        &  86 & Inductive type, structural lemmas \\
\texttt{Constructors.lean} & 112 & All derived operations \\
\texttt{Diff.lean}         & 106 & Differentiation function \\
\texttt{Rewrite.lean}      & 257 & 20 rewrite rules, closure, simplifier \\
\texttt{Theorems.lean}     & 226 & Congruence, $D(\text{const})=0$, $D(e^x)=e^x$ \\
\texttt{Liouville.lean}    & 588 & Chain rules, linearity, cancellation \\
\texttt{Partiality.lean}   &  93 & Non-confluence theorem \\
\end{tabular}

\medskip\noindent
Every theorem in Sections~\ref{sec:diff}--\ref{sec:partiality}
corresponds to a named, compiled Lean definition.  The formalization
is available at \texttt{[repository URL]}.

%% -----------------------------------------------------------------
\section{Discussion}
\label{sec:discussion}
%% -----------------------------------------------------------------

\paragraph{Differentiation without analysis.}
The results of Section~\ref{sec:diff} establish the standard
differentiation rules---chain rules, linearity, specific
derivatives---without invoking limits, real numbers, or
$\varepsilon$-$\delta$ arguments.  The proofs are finite rewrite
chains: purely syntactic, computationally verifiable, and independent
of any model.  This is possible because the EML encoding bakes the
chain rule into the recursive structure of differentiation itself
(equation~\eqref{eq:diff}).

\paragraph{Partiality from a single point.}
The observation that all domain restrictions in elementary mathematics
trace to $\ln 0$ (Section~\ref{sec:partiality}) appears to be new.
In standard presentations, the restrictions $\ln 0$, $1/0$, and
$0^0$ are introduced separately.  In the EML encoding, they collapse:
$1/z = e^{-\ln z}$ and $z^w = e^{w \ln z}$ both route through $\ln z$,
so their singularity at $z = 0$ is inherited from the single constraint
$y \neq 0$ in $\operatorname{eml}(x,y)$.

\paragraph{The tree algebra as a free structure.}
EML trees under the rewrite relation form a free magma---a set with a
binary operation, no associativity, no identity, no inverse.  From this
bare structure, field operations, exponentiation, logarithms, and
differentiation emerge as \emph{patterns}, not axioms.  The algebraic
identities of Section~\ref{sec:algebra} are consequences of the
rewrite rules interacting with the specific tree encodings of $+$,
$\times$, etc.  This suggests viewing the EML tree algebra as the
``free exponential-logarithmic algebra''---the universal structure
generated by one binary operation in which $\exp$ and $\ln$ are total
mutual inverses.

\paragraph{Semantic interpretation.}
The rewrite system is purely syntactic.  Connecting EML trees to
real- or complex-valued functions requires interpreting them in a
model.  The natural candidate---$(\mathbb{R}, +, \times, \exp,
\ln)$---is not a direct model, because $\ln$ is partial on
$\mathbb{R}$ (undefined at $x \le 0$) and multi-valued on
$\mathbb{C}$.  The EML tree algebra is in this sense \emph{richer}
than the reals: it is the structure where $\exp$ and $\ln$ work
everywhere.  Real and complex analysis arise by restricting to
well-formed trees---those whose evaluation never passes $\ln$
through zero.  Formalizing this semantic bridge, including a
well-formedness predicate and a proof that the rewrite rules
preserve it, is natural future work.

\paragraph{Open questions.}
\begin{enumerate}
\item \emph{Confluence.}
  Is the rewrite system confluent on well-formed trees (those whose
  evaluation avoids $\ln 0$)?
  Theorem~\ref{thm:nonconf} shows non-confluence at $1/0$;
  we conjecture this is the only obstruction.
\item \emph{Product and quotient rules.}
  $D(fg) \to^{*} fD(g) + gD(f)$ and $D(f/g) \to^{*}
  (gD(f) - fD(g))/g^2$ should follow from the chain rules and
  linearity but have not yet been formalized.
\item \emph{Completeness of the rewrite system.}
  Do the 20 rules suffice to prove every identity that holds in
  all exponential fields, or are additional rules needed?
\item \emph{Liouville's theorem.}
  Can the classical theorem on elementary antiderivatives be stated
  and proved in the EML tree algebra?  Partial progress exists:
  the Liouville normal form and its derivative structure theorem
  are formalized, but the main induction requires axioms about the
  algebraic independence of $\exp$ and $\ln$.
\end{enumerate}

%% -----------------------------------------------------------------
\bibliographystyle{plain}
\begin{thebibliography}{9}

\bibitem{odrzywolek2026}
Andrzej Odrzywołek.
\newblock All elementary functions from a single operator.
\newblock \emph{arXiv preprint arXiv:2603.21852}, 2026.

\bibitem{lean4}
Leonardo de~Moura and Sebastian Ullrich.
\newblock The {Lean}~4 theorem prover and programming language.
\newblock In \emph{CADE-28}, pages 625--635, 2021.

\bibitem{sheffer1913}
Henry~Maurice Sheffer.
\newblock A set of five independent postulates for {Boolean} algebras.
\newblock \emph{Trans.\ Amer.\ Math.\ Soc.}, 14(4):481--488, 1913.

\bibitem{liouville1835}
Joseph Liouville.
\newblock M{\'e}moire sur l'int{\'e}gration d'une classe de fonctions
  transcendantes.
\newblock \emph{J.\ reine angew.\ Math.}, 13:93--118, 1835.

\bibitem{ritt1948}
Joseph~Fels Ritt.
\newblock \emph{Integration in Finite Terms}.
\newblock Columbia University Press, 1948.

\end{thebibliography}

\end{document}
